\documentclass[10pt,letterpaper]{article}

\usepackage[T1]{fontenc}
\usepackage[utf8]{inputenc}
\usepackage[margin=0.80in,headheight=14pt]{geometry}
\usepackage{amsmath,amssymb,amsthm,mathtools}
\usepackage{newtxtext}
\usepackage{newtxmath}
\DeclareMathAlphabet{\mathcal}{OMS}{cmsy}{m}{n}
\usepackage{bm}
\usepackage{microtype}
\usepackage{booktabs}
\usepackage{array}
\usepackage{enumitem}
\usepackage{xcolor}
\usepackage[most]{tcolorbox}
\usepackage{fancyhdr}
\usepackage[colorlinks=true,linkcolor=blue!65!black,citecolor=blue!65!black,urlcolor=blue!65!black]{hyperref}

\emergencystretch=1.8em

% Pedagogical color palette for Version 3 (Bottom-Up First-Principles Course)
\definecolor{tealborder}{HTML}{0E6251}
\definecolor{softteal}{HTML}{E8F8F5}
\definecolor{crimsonborder}{HTML}{922B21}
\definecolor{softcrimson}{HTML}{FDEDEC}
\definecolor{forestborder}{HTML}{145A32}
\definecolor{softgreen}{HTML}{EAFAF1}
\definecolor{slateborder}{HTML}{2C3E50}
\definecolor{softslate}{HTML}{F2F4F4}

\newtcolorbox{foundationbox}[1]{
  breakable,enhanced,lines before break=2,
  colback=softteal,colframe=tealborder,fonttitle=\bfseries\small,
  title={First-Principles Building Block [BACKGROUND --- Not in Main Paper]: #1},
  left=8pt,right=8pt,top=5pt,bottom=5pt,before skip=7pt,after skip=7pt
}

\newtcolorbox{paperhit}[1]{
  breakable,enhanced,lines before break=2,
  colback=softcrimson,colframe=crimsonborder,fonttitle=\bfseries\small,
  title={Paper Milestone [ALREADY IN JOSA A]: #1},
  left=8pt,right=8pt,top=5pt,bottom=5pt,before skip=7pt,after skip=7pt
}

\newtcolorbox{pencilbox}[1][Pencil-and-Paper Checkpoint]{
  breakable,enhanced,lines before break=2,
  colback=softgreen,colframe=forestborder,fonttitle=\bfseries\small,
  title={#1},
  left=8pt,right=8pt,top=5pt,bottom=5pt,before skip=7pt,after skip=7pt
}

\newtcolorbox{whystep}[1][Why We Take the Next Step]{
  breakable,enhanced,lines before break=2,
  colback=softslate,colframe=slateborder,fonttitle=\bfseries\small,
  title={#1},
  left=8pt,right=8pt,top=5pt,bottom=5pt,before skip=7pt,after skip=7pt
}

\pagestyle{fancy}
\fancyhf{}
\fancyhead[L]{\footnotesize\color{slateborder}\textit{JOSA A (V15 Paper 1) --- Version 3: Bottom-Up First-Principles Course}}
\fancyhead[R]{\footnotesize\color{slateborder}\textit{S. Ayala}}
\fancyfoot[C]{\footnotesize\thepage}
\renewcommand{\headrulewidth}{0.4pt}

\newcommand{\braket}[2]{\langle #1 | #2 \rangle}
\newcommand{\ket}[1]{\lvert #1 \rangle}
\newcommand{\bra}[1]{\langle #1 \rvert}
\newcommand{\norm}[1]{\lVert #1 \rVert}
\newcommand{\abs}[1]{\lvert #1 \rvert}
\newcommand{\dd}{\mathrm{d}}
\newcommand{\inpaper}[1]{\textcolor{crimsonborder}{\textbf{[In JOSA~A: #1]}}}
\newcommand{\bgtag}{\textcolor{tealborder}{\textbf{[Background]}}}

\numberwithin{equation}{section}

\title{\vspace{-1.0cm}\textbf{\Large Version 3 --- From First Principles With Pencil and Paper (JOSA A, Paper 1)}\\[0.35em]
\large\textit{A Bottom-Up Self-Study Course Starting From Maxwell's Equations and Gaussian Integrals,\\
Building Step by Step to Every Equation (1)--(26) of ``Continuous Hilbert frame theory, paraxial propagation invariance,\\
and exact Laguerre--Gaussian expansion of finite-width perfect optical vortex beams''}}
\author{Sebasti\'an Ayala\thanks{\texttt{ayala.phys.ds@gmail.com}; ORCID: \href{https://orcid.org/0009-0003-7348-3883}{0009-0003-7348-3883}}}
\date{Version 3 (First-Principles Self-Study Edition) --- November 24, 2026}

\begin{document}
\maketitle
\thispagestyle{plain}

\begin{abstract}
\noindent Unlike Versions~1 and~2 (which follow the published paper's top-down section order), this \textbf{Version~3} course is written for someone \textbf{completely new to structured light and frame theory} who wants to pick up a pencil and paper and derive the entire JOSA~A manuscript from scratch. We start from the simplest universal equations of undergraduate physics and calculus---Maxwell's equations, the 1D Gaussian integral $\int_{-\infty}^\infty e^{-ax^2}\dd x = \sqrt{\pi/a}$, polar Fourier transforms, and a 3-vector toy frame in $\mathbb{R}^2$---and build organically upward. Every foundational warm-up is clearly marked \bgtag{}, and the exact moment our bottom-up derivation arrives at a result published in \texttt{main\_paper1\_josaa.pdf}, a crimson \emph{Paper Milestone} box marks \inpaper{Eq.~(k), p.~m}.
\end{abstract}

\tableofcontents
\vspace{0.5em}

% ==============================================================================
\section{Chapter 1: From Maxwell's Equations to Paraxial Light Beams \bgtag{}}
\label{sec:c1}

\begin{whystep}[Our Starting Point]
All of Paper~1 lives inside a single physical arena: monochromatic laser beams propagating mostly along one axis ($z$). Before talking about "vortex beams" or "frames," let us derive the governing wave equation directly from Maxwell's equations so you see every approximation with your own eyes.
\end{whystep}

\subsection{Step 1.1: Maxwell's Equations $\to$ Scalar Helmholtz Equation}

In empty space (no free charges or currents, permittivity $\varepsilon_0$, permeability $\mu_0$, speed of light $c = 1/\sqrt{\mu_0\varepsilon_0}$), Maxwell's two curl equations and Gauss's law are:
\begin{equation}
  \nabla\times\mathbf{E} = -\frac{\partial\mathbf{B}}{\partial t}, \qquad
  \nabla\times\mathbf{B} = \frac{1}{c^2}\frac{\partial\mathbf{E}}{\partial t}, \qquad
  \nabla\cdot\mathbf{E} = 0.
\end{equation}
Take the curl of the first equation, use the vector calculus identity $\nabla\times(\nabla\times\mathbf{E}) = \nabla(\nabla\cdot\mathbf{E}) - \nabla^2\mathbf{E} = -\nabla^2\mathbf{E}$, and substitute the second equation:
\begin{equation}
  -\nabla^2\mathbf{E} = -\frac{\partial}{\partial t}(\nabla\times\mathbf{B}) = -\frac{1}{c^2}\frac{\partial^2\mathbf{E}}{\partial t^2}
  \implies \left(\nabla^2 - \frac{1}{c^2}\frac{\partial^2}{\partial t^2}\right)\mathbf{E}(\mathbf{r},z,t) = 0.
\end{equation}
For a single-color (monochromatic) laser beam of angular frequency $\omega = c k_0$ ($k_0 = 2\pi/\lambda$) polarized uniformly along $\hat{\mathbf{x}}$, write $\mathbf{E}(\mathbf{r},z,t) = \hat{\mathbf{x}}\,\operatorname{Re}[E(\mathbf{r},z)e^{-i\omega t}]$. Since $\partial_t^2 e^{-i\omega t} = -\omega^2 e^{-i\omega t} = -c^2 k_0^2 e^{-i\omega t}$, the spatial complex amplitude $E(\mathbf{r},z)$ obeys the \textbf{scalar Helmholtz equation}:
\begin{equation}
  \left(\nabla_\perp^2 + \frac{\partial^2}{\partial z^2} + k_0^2\right)E(\mathbf{r},z) = 0, \qquad \nabla_\perp^2 \equiv \frac{\partial^2}{\partial x^2} + \frac{\partial^2}{\partial y^2}.
  \label{eq:c1_helmholtz}
\end{equation}

\subsection{Step 1.2: Slowly Varying Envelope $\to$ Paraxial Wave Equation}

A laser beam travels primarily along $+z$, so its fastest oscillation is the plane-wave carrier $e^{+ik_0 z}$. Factor it out by writing $E(\mathbf{r},z) = \psi(\mathbf{r},z)\,e^{ik_0 z}$ and differentiate twice with respect to $z$:
\begin{align}
  \frac{\partial E}{\partial z} &= \left(\frac{\partial\psi}{\partial z} + i k_0\psi\right)e^{ik_0 z}, \\
  \frac{\partial^2 E}{\partial z^2} &= \left(\frac{\partial^2\psi}{\partial z^2} + 2ik_0\frac{\partial\psi}{\partial z} - k_0^2\psi\right)e^{ik_0 z}.
\end{align}
Insert into Eq.~\eqref{eq:c1_helmholtz}: the $-k_0^2\psi$ and $+k_0^2\psi$ terms cancel \emph{identically}, leaving $(\nabla_\perp^2 + 2ik_0\partial_z + \partial_z^2)\psi = 0$. Now make the \textbf{paraxial approximation}: if the beam spreads slowly compared to one wavelength, $\abs{\partial_z^2\psi} \ll 2k_0\abs{\partial_z\psi}$, so we drop $\partial_z^2\psi$ and obtain the \textbf{paraxial Helmholtz equation}:
\begin{equation}
  \left(2ik_0\frac{\partial}{\partial z} + \nabla_\perp^2\right)\psi(\mathbf{r},z) = 0
  \iff i\frac{\partial\psi}{\partial z} = -\frac{1}{2k_0}\nabla_\perp^2\psi.
  \label{eq:c1_paraxial}
\end{equation}
Notice the astonishing analogy: Eq.~\eqref{eq:c1_paraxial} has the \emph{exact mathematical form} of the 2D free-particle Schr\"odinger equation ($i\hbar\partial_t\psi = -\frac{\hbar^2}{2m}\nabla_\perp^2\psi$) with axial distance $z$ playing the role of time $t$!

\subsection{Step 1.3: Solving Eq.~\eqref{eq:c1_paraxial} Instantly via 2D Fourier Transform}

Define the 2D transverse Fourier transform of any square-integrable beam $\psi(\mathbf{r},z) \in L^2(\mathbb{R}^2)$:
\begin{equation}
  \widetilde{\psi}(\mathbf{k},z) \equiv \int_{\mathbb{R}^2} \psi(\mathbf{r},z) e^{-i\mathbf{k}\cdot\mathbf{r}}\,\dd^2\mathbf{r}, \qquad
  \psi(\mathbf{r},z) = \int_{\mathbb{R}^2} \widetilde{\psi}(\mathbf{k},z) e^{+i\mathbf{k}\cdot\mathbf{r}}\,\frac{\dd^2\mathbf{k}}{(2\pi)^2}.
  \label{eq:c1_ft}
\end{equation}
Under integration by parts (since $\psi \to 0$ at $\infty$), $\partial_x \mapsto +ik_x$ and $\nabla_\perp^2 \mapsto -(k_x^2 + k_y^2) = -k_\perp^2$. Thus the PDE~\eqref{eq:c1_paraxial} collapses into an elementary 1st-order ODE for each spatial frequency $\mathbf{k}$:
\begin{equation}
  2ik_0\frac{\partial\widetilde{\psi}}{\partial z} - k_\perp^2\widetilde{\psi} = 0
  \implies \widetilde{\psi}(\mathbf{k},z) = \exp\!\left(-\frac{iz k_\perp^2}{2k_0}\right)\widetilde{\psi}(\mathbf{k},0).
  \label{eq:c1_fresnel}
\end{equation}
We write $\psi(\cdot, z) = \mathcal{U}_z\psi(\cdot, 0)$. Because $\abs{e^{-iz k_\perp^2/(2k_0)}} = 1$, Plancherel's theorem ($\int\abs{\psi}^2\dd^2\mathbf{r} = \int\abs{\widetilde{\psi}}^2\frac{\dd^2\mathbf{k}}{(2\pi)^2}$) immediately shows $\norm{\mathcal{U}_z\psi}_{L^2}^2 = \norm{\psi}_{L^2}^2$: free-space propagation conserves total optical power (\textbf{unitarity}: $\mathcal{U}_z^\dagger = \mathcal{U}_z^{-1} = \mathcal{U}_{-z}$)!

\begin{foundationbox}{The Simplest Beam Solution: Gaussian Beam, Curvature $R_c(z)$, and Gouy Phase $\vartheta(z)$}
Start at $z = 0$ with a simple Gaussian spot of waist $w_0$: $\psi_{0,0}(r,0) = \sqrt{\frac{2}{\pi}}\frac{1}{w_0} e^{-r^2/w_0^2}$ (check $\int_0^\infty 2\pi r\abs{\psi_{0,0}}^2\dd r = 1$!). Its 2D Fourier transform (using Chapter~2's Gaussian integral) is $\widetilde{\psi}_{0,0}(\mathbf{k},0) = w_0\sqrt{2\pi}\,e^{-w_0^2 k_\perp^2/4}$. Multiply by Eq.~\eqref{eq:c1_fresnel}:
\[
  -\frac{w_0^2 k_\perp^2}{4} - \frac{iz k_\perp^2}{2k_0} = -\frac{k_\perp^2}{4}\,w_0^2\!\left(1 + i\frac{z}{z_R}\right), \qquad z_R \equiv \frac{k_0 w_0^2}{2}\quad\text{(Rayleigh range)}.
\]
Inverting the Gaussian Fourier transform replaces $w_0^2 \mapsto w_0^2(1 + iz/z_R) = w_0^2\sigma(z)e^{i\vartheta(z)}$ where $\sigma(z) = \sqrt{1+(z/z_R)^2}$ and $\vartheta(z) = \arctan(z/z_R)$. Rationalizing $1/[w_0^2(1+iz/z_R)] = \frac{1}{w_0^2\sigma^2(z)} - \frac{ik_0}{2R_c(z)}$ (with $R_c(z) = z[1+(z_R/z)^2]$) yields the famous propagating Gaussian beam:
\[
  \psi_{0,0}(r,z) = \frac{1}{\sigma(z)}\,\psi_{0,0}\!\left(\frac{r}{\sigma(z)},0\right)\exp\!\left[+\frac{ik_0 r^2}{2R_c(z)} - i\vartheta(z)\right].
\]
Notice the extra on-axis phase $-\vartheta(z) = -\arctan(z/z_R)$: that is the \textbf{Gouy phase}, acquired purely because a focused beam is transversely confined!
\end{foundationbox}

% ==============================================================================
\section{Chapter 2: Essential Pencil-and-Paper Math Toolkit (Derived From Zero!) \bgtag{}}
\label{sec:c2}

\begin{whystep}[Zero Black Boxes]
Every single calculation in Paper~1 relies on just \textbf{four} elementary mathematical tools: (1) Gaussian integrals, (2) Jacobi--Anger angular integration, (3) Weber's Gaussian--Bessel integral, and (4) Hankel's large-$x$ expansion of $I_n(x)$. Let us derive ALL FOUR from scratch right now so nothing ever feels like magic!
\end{whystep}

\subsection{Step 2.1: The 1D Gaussian Integral and Completing the Square}

Let $I \equiv \int_{-\infty}^\infty e^{-ax^2}\,\dd x$ ($a > 0$). Square it and switch to polar coordinates $(x,y) = (r\cos\phi, r\sin\phi)$ ($\dd x\,\dd y = r\,\dd r\,\dd\phi$):
\begin{equation}
  I^2 = \int_{-\infty}^\infty\!\!\int_{-\infty}^\infty e^{-a(x^2+y^2)}\,\dd x\,\dd y = \int_0^{2\pi}\dd\phi \int_0^\infty r\,e^{-ar^2}\,\dd r = 2\pi\left[-\frac{e^{-ar^2}}{2a}\right]_0^\infty = \frac{\pi}{a}
  \implies \int_{-\infty}^\infty e^{-ax^2}\,\dd x = \sqrt{\frac{\pi}{a}}.
  \label{eq:c2_gauss}
\end{equation}
Adding a linear term $bx$ (even complex $b \in \mathbb{C}$, $\operatorname{Re} a > 0$) works by completing the square $-ax^2 + bx = -a(x - \frac{b}{2a})^2 + \frac{b^2}{4a}$:
\begin{equation}
  \int_{-\infty}^\infty e^{-ax^2 + bx}\,\dd x = e^{b^2/(4a)}\int_{-\infty}^\infty e^{-a(x - b/(2a))^2}\,\dd x = \sqrt{\frac{\pi}{a}}\,e^{b^2/(4a)}.
  \label{eq:c2_gauss_shift}
\end{equation}
Multiplying two copies in $x$ and $y$ immediately gives the 2D Gaussian Fourier identity:
\begin{equation}
  \int_{\mathbb{R}^2} e^{-p\abs{\mathbf{x}}^2 + i\mathbf{q}\cdot\mathbf{x}}\,\dd^2\mathbf{x} = \frac{\pi}{p}\exp\!\left(-\frac{\abs{\mathbf{q}}^2}{4p}\right) \qquad (\operatorname{Re} p > 0).
  \label{eq:c2_gauss_2d}
\end{equation}

\subsection{Step 2.2: Polar Coordinates $\to$ Bessel Functions $J_n(z)$ and Modified Bessel Functions $I_n(z)$}

Why do Bessel functions appear everywhere in circular beams? Because whenever you Fourier-transform a wave carrying circular phase $e^{il\phi}$, you encounter the angular integral $\int_0^{2\pi} e^{il\phi} e^{-i z\cos(\phi-\phi_k)}\,\dd\phi$!
Define the integer-order Bessel function $J_n(z)$ and modified Bessel function $I_n(z)$ directly as the $n$-th Fourier coefficients of the periodic functions $e^{+iz\cos\psi}$ and $e^{+z\cos\psi}$:
\begin{align}
  e^{\pm i z\cos\psi} &= \sum_{n=-\infty}^\infty (\pm i)^n J_n(z) e^{in\psi}
  \iff J_n(z) = \frac{(\mp i)^n}{2\pi}\int_0^{2\pi} e^{\pm iz\cos\psi} e^{-in\psi}\,\dd\psi \quad\text{(Jacobi--Anger)},
  \label{eq:c2_ja} \\
  e^{z\cos\psi} &= \sum_{n=-\infty}^\infty I_n(z) e^{in\psi}
  \iff I_n(z) = \frac{1}{2\pi}\int_0^{2\pi} e^{z\cos\psi} e^{-in\psi}\,\dd\psi = \frac{1}{\pi}\int_0^\pi e^{z\cos\psi}\cos(n\psi)\,\dd\psi.
  \label{eq:c2_In}
\end{align}
Look how many "textbook facts" drop out of Eqs.~\eqref{eq:c2_ja}--\eqref{eq:c2_In} in ONE line of mental math:
\begin{itemize}[leftmargin=1.4em,itemsep=1pt,topsep=1pt]
  \item Replacing $z \mapsto iz$ in Eq.~\eqref{eq:c2_ja} (with lower sign $-i$) turns $e^{-i(iz)\cos\psi} = e^{z\cos\psi}$, giving $(-i)^n J_n(iz) = I_n(z)$, i.e.\ $J_n(iz) = i^n I_n(z)$ and $I_n(-iz) = (-i)^n J_n(z)$!
  \item Setting $\psi \mapsto -\psi$ in Eq.~\eqref{eq:c2_In} gives $I_{-n}(z) = I_n(z) = I_{\abs{n}}(z)$ (even in $n$!), while in Eq.~\eqref{eq:c2_ja} it gives $J_{-n}(z) = (-1)^n J_n(z)$!
  \item Setting $\psi \mapsto \psi + \pi$ in Eq.~\eqref{eq:c2_In} gives $I_n(-z) = (-1)^n I_n(z)$!
  \item Multiplying Eq.~\eqref{eq:c2_In} by $e^{-z}$ gives $e^{-z} I_n(z) = \frac{1}{\pi}\int_0^\pi e^{-z(1-\cos\psi)}\cos(n\psi)\,\dd\psi$, which is \emph{strictly less than $1$} for every $z > 0$ because $e^{-z(1-\cos\psi)} < 1$ on $(0,\pi)$!
\end{itemize}

\subsection{Step 2.3: Deriving Weber's Master Integral From Scratch!}

Put $\mathbf{q} = \boldsymbol{\alpha} - \boldsymbol{\beta}$ in Eq.~\eqref{eq:c2_gauss_2d}, where $\abs{\boldsymbol{\alpha}} = \alpha$, $\abs{\boldsymbol{\beta}} = \beta$, and $\theta_\alpha - \theta_\beta$ is the angle between them:
\begin{itemize}[leftmargin=1.4em,itemsep=1pt,topsep=1pt]
  \item \textbf{Right-hand side:} Expand $\abs{\boldsymbol{\alpha}-\boldsymbol{\beta}}^2 = \alpha^2 + \beta^2 - 2\alpha\beta\cos(\theta_\alpha - \theta_\beta)$ and apply Eq.~\eqref{eq:c2_In}:
  \[
    \frac{\pi}{p}\exp\!\left(-\frac{\alpha^2+\beta^2}{4p}\right) e^{\frac{\alpha\beta}{2p}\cos(\theta_\alpha-\theta_\beta)}
    = \frac{\pi}{p}\exp\!\left(-\frac{\alpha^2+\beta^2}{4p}\right)\sum_{n=-\infty}^\infty I_{\abs{n}}\!\left(\frac{\alpha\beta}{2p}\right) e^{in(\theta_\alpha-\theta_\beta)}.
  \]
  \item \textbf{Left-hand side:} Write $\mathbf{x} = (x,\varphi)$ ($\dd^2\mathbf{x} = x\,\dd x\,\dd\varphi$), expand $e^{i\alpha x\cos(\varphi-\theta_\alpha)}$ and $e^{-i\beta x\cos(\varphi-\theta_\beta)}$ via Eq.~\eqref{eq:c2_ja}, and integrate $\int_0^{2\pi} e^{i(n-m)\varphi}\,\dd\varphi = 2\pi\delta_{nm}$:
  \[
    \text{LHS} = 2\pi\sum_{n=-\infty}^\infty e^{in(\theta_\alpha-\theta_\beta)}\int_0^\infty x\,e^{-px^2} J_{\abs{n}}(\alpha x) J_{\abs{n}}(\beta x)\,\dd x.
  \]
\end{itemize}
Equating the $n$-th Fourier coefficient on both sides gives \textbf{Weber's second exponential integral} ($\operatorname{Re} p > 0$, $n = \abs{l} \in \mathbb{N}_0$):
\begin{equation}
  \int_0^\infty x\,e^{-px^2} J_{\abs{l}}(\alpha x) J_{\abs{l}}(\beta x)\,\dd x = \frac{1}{2p}\exp\!\left(-\frac{\alpha^2+\beta^2}{4p}\right) I_{\abs{l}}\!\left(\frac{\alpha\beta}{2p}\right).
  \label{eq:c2_weberJ}
\end{equation}
And by analytic continuation $\alpha \to ia, \beta \to ib$ (using $J_{\abs{l}}(iy) = i^{\abs{l}}I_{\abs{l}}(y)$ and $I_{\abs{l}}(-y) = (-1)^{\abs{l}}I_{\abs{l}}(y)$), we immediately get its modified partner:
\begin{equation}
  \int_0^\infty x\,e^{-px^2} I_{\abs{l}}(ax) I_{\abs{l}}(bx)\,\dd x = \frac{1}{2p}\exp\!\left(\frac{a^2+b^2}{4p}\right) I_{\abs{l}}\!\left(\frac{ab}{2p}\right).
  \label{eq:c2_weberI}
\end{equation}

\subsection{Step 2.4: Deriving the Large-$x$ Behavior of $I_n(x)$ (Hankel Expansion) From Scratch!}

What does $I_n(x) = \frac{1}{\pi}\int_0^\pi e^{x\cos\psi}\cos(n\psi)\,\dd\psi$ look like when $x \gg 1$? Near $\psi = 0$, $\cos\psi \approx 1 - \psi^2/2$ peaks sharply, while everywhere else $e^{x\cos\psi}$ is exponentially tiny! Applying Eq.~\eqref{eq:c2_gauss}:
\begin{equation}
  I_n(x) \approx \frac{e^x}{\pi}\int_0^\infty e^{-x\psi^2/2}\cdot 1\,\dd\psi = \frac{e^x}{\pi}\cdot\frac{1}{2}\sqrt{\frac{2\pi}{x}} = \frac{e^x}{\sqrt{2\pi x}}.
\end{equation}
To get the systematic $1/x$ corrections, differentiate Eq.~\eqref{eq:c2_In} under the integral sign to verify $I_n(x)$ obeys $x^2 y'' + x y' - (x^2+n^2)y = 0$, substitute $y(x) = \frac{e^x}{\sqrt{2\pi x}}\sum_{k=0}^\infty c_k x^{-k}$ ($c_0 = 1$), and equate powers of $x^{-k}$:
\begin{equation}
  c_{k+1} = -\frac{4n^2 - (2k+1)^2}{8(k+1)}\,c_k
  \implies I_n(x) \sim \frac{e^x}{\sqrt{2\pi x}}\left[1 - \frac{\mu}{8x} + \frac{\mu(\mu-8)}{128 x^2} - \cdots\right], \qquad \mu \equiv 4n^2 - 1.
  \label{eq:c2_hankel}
\end{equation}

\begin{pencilbox}[Quick Test of Your Toolkit Before Entering Paper 1]
Check Eq.~\eqref{eq:c2_hankel} at $n = 1$ ($\mu = 3$) and $x = 16$: $(2\pi\cdot 16)^{-1/2} = 0.099736$; $1 - 3/128 = 0.976563$; product $= 0.097398$. The exact value of $e^{-16}I_1(16)$ is $0.097350$ (accurate to $0.05\%$ from just two terms!).
\end{pencilbox}

% ==============================================================================
\section{Chapter 3: Building a Physical Perfect Vortex Beam From a Lens}
\label{sec:c3}

\begin{whystep}[Why Invent a "Perfect" Vortex Beam?]
A vortex beam carries twisted phase $e^{il\phi}$ (orbital angular momentum $l\hbar$ per photon). Standard Laguerre--Gaussian ($\mathrm{LG}_{0,l}$) laser modes have radial intensity $\propto r^{2\abs{l}} e^{-2r^2/w_0^2}$, whose peak sits at $r_{\mathrm{peak}} = w_0\sqrt{\abs{l}/2}$---so higher $l$ forces a bigger ring! That breaks coupling into fixed-radius ring-core optical fibers. Can we build a \emph{physical} light ring whose radius $R$ stays fixed for every $l$?
\end{whystep}

\subsection{Step 3.1: From an Ideal Delta Ring to a $2f$ Fourier Lens}

An idealized mathematical ring $\delta(r-R)e^{il\phi}$ has infinite energy ($\int \delta(r-R)^2 r\,\dd r = \infty$). Notice, however, what its 2D Fourier transform would be: a pure Bessel beam $J_l(k_r\rho)e^{il\theta}$ (which itself carries infinite energy because $J_l$ decays slowly as $\rho^{-1/2}$). To make it \emph{physical} (finite energy), we multiply the Bessel beam by a Gaussian laser envelope $e^{-\rho^2/w_g^2}$ of finite waist $w_g$:
\begin{equation}
  E_{\mathrm{BG}}(\rho,\theta) = E_0 J_l(k_r\rho) e^{-\rho^2/w_g^2} e^{il\theta},
\end{equation}
and place it at the front focal plane of a thin spherical lens of focal length $f$. In paraxial Fourier optics, a $2f$ lens computes the exact 2D optical Fourier transform at its back focal plane:
\begin{equation}
  E_{\mathrm{out}}(r,\phi) = \frac{e^{2ik_0 f}}{i\lambda f}\int_0^\infty \rho\,\dd\rho\,E_0 J_l(k_r\rho)e^{-\rho^2/w_g^2}\int_0^{2\pi}\dd\theta\,e^{il\theta}\exp\!\left[-\frac{ik_0}{f}r\rho\cos(\theta-\phi)\right].
\end{equation}
Apply our Chapter~2 Jacobi--Anger identity~\eqref{eq:c2_ja} to the $\theta$-integral (picking out $n = -l$, giving $2\pi(-i)^{\abs{l}}J_{\abs{l}}(\frac{k_0 r}{f}\rho)e^{il\phi}$), then apply Weber's integral~\eqref{eq:c2_weberJ} with $p = 1/w_g^2$, $\alpha = k_r$, $\beta = k_0 r/f$:
\begin{equation}
  E_{\mathrm{out}}(r,\phi) = (-i)^{l+1} e^{2ik_0 f}\left(\frac{k_0 E_0 w_g^2}{2f}\right) e^{il\phi}\exp\!\left[-\frac{w_g^2 k_r^2}{4} - \frac{w_g^2 k_0^2}{4f^2}r^2\right] I_{\abs{l}}\!\left(\frac{w_g^2 k_r k_0}{2f}r\right).
\end{equation}
Now define two physical length scales that make the exponent clean: the \textbf{ring radius} $R \equiv \frac{k_r f}{k_0}$ and the \textbf{ring half-width} $w \equiv \frac{2f}{k_0 w_g}$. Check the algebra: $\frac{w_g^2 k_0^2}{4f^2} = \frac{1}{w^2}$, $\frac{w_g^2 k_r^2}{4} = \frac{R^2}{w^2}$, and $\frac{w_g^2 k_r k_0}{2f}r = \frac{2Rr}{w^2}$! Choosing the input amplitude $E_0 = \frac{k_0}{f\sqrt{2\pi}} = \frac{1}{w w_g}\sqrt{\frac{2}{\pi}}$ fixes the overall prefactor to $C_0 \equiv \frac{1}{w^2}\sqrt{\frac{2}{\pi}}$ (whose meaning becomes clear in Chapter~4!).

\begin{paperhit}{JOSA A Eq.~(1) (p.~2) and Eq.~(4) (p.~3)}
We have just derived the central beam family of Paper~1 and its laboratory $2f$ realization from scratch:
\[
  \Phi_{l,R}(r,\phi) = \frac{1}{w^2}\sqrt{\frac{2}{\pi}}\,e^{il\phi}\exp\!\left(-\frac{r^2+R^2}{w^2}\right) I_{\abs{l}}\!\left(\frac{2Rr}{w^2}\right), \qquad
  R = \frac{k_r f}{k_0}, \quad w = \frac{2f}{k_0 w_g}, \quad E_0 = \frac{k_0}{f\sqrt{2\pi}}.
\]
\end{paperhit}

\subsection{Step 3.2: Total Optical Power (Exact $L^2$ Norm)}

How much power is in $\Phi_{l,R}$? Integrate $\abs{\Phi_{l,R}}^2$ over the 2D plane using $\int_0^{2\pi}\dd\phi = 2\pi$ and our modified Weber integral~\eqref{eq:c2_weberI} with $p = 2/w^2$, $a = b = 2R/w^2$:
\begin{equation}
  \norm{\Phi_{l,R}}_{L^2}^2 = 2\pi C_0^2 e^{-2R^2/w^2}\cdot\frac{w^2}{4} e^{+R^2/w^2} I_{\abs{l}}\!\left(\frac{R^2}{w^2}\right) = \frac{1}{w^2} e^{-R^2/w^2} I_{\abs{l}}\!\left(\frac{R^2}{w^2}\right) < \frac{1}{w^2},
\end{equation}
where the strict inequality $< w^{-2}$ follows directly from Step~2.2 ($e^{-x} I_n(x) < 1$ for $x > 0$)!

\begin{paperhit}{JOSA A Eq.~(2) (p.~2)}
Exact finite energy bound of every physical PVB: $\norm{\Phi_{l,R}}_{L^2}^2 = \frac{1}{w^2} e^{-R^2/w^2} I_{\abs{l}}(R^2/w^2) < \frac{1}{w^2}$.
\end{paperhit}

\subsection{Step 3.3: Is the Ring Radius Truly Independent of $l$?}

Plug our leading Hankel formula $I_{\abs{l}}(x) \approx e^x/\sqrt{2\pi x}$ into $\Phi_{l,R}$: the exponents combine into $-(r-R)^2/w^2$, revealing a Gaussian ring of thickness $w$ centered near $r = R$! To find the \emph{exact} intensity peak, set $\frac{\dd}{\dd r}\ln\abs{\Phi_{l,R}(r)}^2 = 0$:
\begin{equation}
  -\frac{4r}{w^2} + \frac{4R}{w^2}\frac{I'_{\abs{l}}(x)}{I_{\abs{l}}(x)} = 0 \iff \frac{r_{\mathrm{peak}}(l)}{R} = \frac{I'_{\abs{l}}(x_{\mathrm{peak}})}{I_{\abs{l}}(x_{\mathrm{peak}})}, \qquad x_{\mathrm{peak}} = \frac{2R\,r_{\mathrm{peak}}(l)}{w^2}.
\end{equation}
Using Eq.~\eqref{eq:c2_hankel}, $\frac{\dd}{\dd x}\ln I_{\abs{l}}(x) = 1 - \frac{1}{2x} + \frac{4l^2-1}{8x^2} + \mathcal{O}(x^{-3})$. Let $\delta \equiv (w/R)^2 \ll 1$ and try $r_{\mathrm{peak}}/R = 1 + c_1\delta + c_2\delta^2 + \mathcal{O}(\delta^3)$, so $x_{\mathrm{peak}}^{-1} = \frac{\delta}{2}(1 - c_1\delta) + \mathcal{O}(\delta^3)$:
\begin{equation}
  1 + c_1\delta + c_2\delta^2 = 1 - \frac{1}{2}\left(\frac{\delta}{2} - \frac{c_1}{2}\delta^2\right) + \frac{4l^2-1}{8}\left(\frac{\delta^2}{4}\right) + \mathcal{O}(\delta^3)
  \implies c_1 = -\frac{1}{4}, \quad c_2 = \frac{4l^2 - 3}{32}.
\end{equation}
Notice that $c_1 = -1/4$ shifts EVERY ring slightly inward by $-w^2/(4R)$, regardless of $l$. Subtracting $r_{\mathrm{peak}}(0)$ leaves the pure $l$-dependent residual shift:

\begin{paperhit}{JOSA A Eq.~(3) (p.~3)}
The thin-ring physical PVB is NOT mathematically $l$-independent, but its residual shift is suppressed to \emph{fourth order} in $w/R$:
\[
  \Delta r_{\mathrm{peak}}(l) \equiv r_{\mathrm{peak}}(l) - r_{\mathrm{peak}}(0) \approx \frac{l^2}{8}\left(\frac{w}{R}\right)^4 R.
\]
At $w/R = 0.1$ and $l = 4$, this is only $2.0\times 10^{-4} R$ ($0.02\%$ of the ring radius!).
\end{paperhit}

% ==============================================================================
\section{Chapter 4: From a 3-Vector Toy Frame to the Continuous PVB Frame Operator}
\label{sec:c4}

\begin{whystep}[What Is a "Frame" If You Have Never Seen One?]
Everyone knows an \textbf{orthonormal basis} (like $\hat{\mathbf{x}}, \hat{\mathbf{y}}$ in $\mathbb{R}^2$): every vector expands uniquely as $\mathbf{v} = \braket{\hat{\mathbf{x}}}{\mathbf{v}}\hat{\mathbf{x}} + \braket{\hat{\mathbf{y}}}{\mathbf{v}}\hat{\mathbf{y}}$. But concentric physical PVB rings of finite thickness $w$ \emph{overlap}! So they CANNOT be orthogonal. What replaces an orthonormal basis when vectors overlap? A \textbf{Hilbert frame}! Let's build the simplest frame on earth first.
\end{whystep}

\begin{foundationbox}{Warm-Up Toy Example: The Mercedes-Benz Frame in $\mathbb{R}^2$}
Take three unit vectors at $120^\circ$ in the plane $\mathbb{R}^2$:
\[
  \mathbf{e}_1 = (1, 0), \qquad \mathbf{e}_2 = \left(-\tfrac{1}{2}, \tfrac{\sqrt{3}}{2}\right), \qquad \mathbf{e}_3 = \left(-\tfrac{1}{2}, -\tfrac{\sqrt{3}}{2}\right).
\]
They are \emph{overcomplete} (3 vectors in a 2D plane) and \emph{non-orthogonal} ($\mathbf{e}_1\cdot\mathbf{e}_2 = -1/2 \neq 0$). Now define the \textbf{frame operator} $S\mathbf{v} \equiv \sum_{k=1}^3 (\mathbf{e}_k\cdot\mathbf{v})\mathbf{e}_k$ for any $\mathbf{v} = (x,y)$ and compute the sum explicitly:
\begin{align*}
  \sum_{k=1}^3 (\mathbf{e}_k\cdot\mathbf{v})^2 &= x^2 + \left(-\frac{x}{2} + \frac{\sqrt{3}}{2}y\right)^2 + \left(-\frac{x}{2} - \frac{\sqrt{3}}{2}y\right)^2 = x^2 + \frac{x^2}{2} + \frac{3y^2}{2} = \frac{3}{2}(x^2 + y^2) = \frac{3}{2}\norm{\mathbf{v}}^2!
\end{align*}
Therefore $S = \frac{3}{2}\mathbb{I}$! Even though the $\mathbf{e}_k$ are NOT orthogonal, you still reconstruct $\mathbf{v}$ \emph{exactly} via $\mathbf{v} = \sum_{k=1}^3 (\mathbf{e}_k\cdot\mathbf{v})(S^{-1}\mathbf{e}_k) = \frac{2}{3}\sum_{k=1}^3 (\mathbf{e}_k\cdot\mathbf{v})\mathbf{e}_k$! In general, a set of vectors is a \textbf{frame} with bounds $0 < A \le B < \infty$ if $A\norm{\mathbf{v}}^2 \le \braket{\mathbf{v}}{S\mathbf{v}} \le B\norm{\mathbf{v}}^2$ for all $\mathbf{v}$---which guarantees $S$ is stably invertible with condition number $\kappa(S) = B/A$!
\end{foundationbox}

\subsection{Step 4.1: Overlap of Two Physical PVBs}

Now return to our continuous family $\{\Phi_{l,R}\}_{l\in\mathbb{Z},\, R>0}$. Compute the inner product $\braket{\Phi_{l,R}}{\Phi_{l',R'}}$: the $\phi$-integral yields $2\pi\delta_{l,l'}$, and Weber's integral~\eqref{eq:c2_weberI} with $p = 2/w^2$, $a = 2R/w^2$, $b = 2R'/w^2$ evaluates the radial integral immediately! Dividing by $\norm{\Phi_{l,R}}\norm{\Phi_{l,R'}}$ gives the normalized overlap $\gamma_l(R,R')$:

\begin{paperhit}{JOSA A Eqs.~(5) and (6) (pp.~3--4)}
\[
  \braket{\Phi_{l,R}}{\Phi_{l',R'}} = \delta_{l,l'}\frac{1}{w^2}e^{-(R^2+R'^2)/(2w^2)} I_{\abs{l}}\!\left(\frac{RR'}{w^2}\right), \qquad
  \gamma_l(R,R') = \frac{I_{\abs{l}}(RR'/w^2)}{\sqrt{I_{\abs{l}}(R^2/w^2)\,I_{\abs{l}}(R'^2/w^2)}} \approx e^{-(R-R')^2/(2w^2)}.
\]
\end{paperhit}

Different OAM values $l \neq l'$ are strictly orthogonal, while concentric rings of the same $l$ overlap with a Gaussian tail of width $\sqrt{2}\,w$---meaning $\{\Phi_{l,R}\}$ is an \emph{overcomplete} radial family!

\subsection{Step 4.2: Building the Continuous PVB Frame Operator $S_w$}

Mimicking our Mercedes-Benz toy sum $S\mathbf{v} = \sum_k \braket{\mathbf{e}_k}{\mathbf{v}}\mathbf{e}_k$, sum over discrete $l \in \mathbb{Z}$ and integrate over continuous radii $R \in (0,\infty)$ with the natural polar area weight $R\,\dd R$:
\begin{equation}
  S_w\psi \equiv \sum_{l\in\mathbb{Z}}\int_0^\infty R\,\dd R\,\braket{\Phi_{l,R}}{\psi}\,\Phi_{l,R}, \qquad
  \braket{\psi}{S_w\psi} = \sum_{l\in\mathbb{Z}}\int_0^\infty R\,\dd R\,\abs{\braket{\Phi_{l,R}}{\psi}}^2 \equiv \norm{T\psi}_{L^2(\Omega,\mu)}^2.
\end{equation}
How do we evaluate $S_w\psi$ cleanly? Go to Fourier space! First compute $\widetilde{\Phi}_{l,R}(\mathbf{k})$ using Jacobi--Anger~\eqref{eq:c2_ja} and Weber~\eqref{eq:c2_weberI} (analytically continuing $b \to -ik_\perp$):
\begin{equation}
  \widetilde{\Phi}_{l,R}(k_\perp,\phi_k) = 2\pi C_0 (-i)^{\abs{l}} e^{il\phi_k} e^{-R^2/w^2} \cdot \frac{w^2}{2} e^{R^2/w^2 - w^2 k_\perp^2/4} J_{\abs{l}}(Rk_\perp) = \sqrt{2\pi}\,(-i)^{\abs{l}} e^{il\phi_k} e^{-w^2 k_\perp^2/4} J_{\abs{l}}(Rk_\perp).
\end{equation}
Look how the $C_0 = \frac{1}{w^2}\sqrt{\frac{2}{\pi}}$ factor was chosen: $2\pi C_0(w^2/2) = \sqrt{2\pi}$ cancels $w$ out of the prefactor completely! Now expand any beam's Fourier transform into angular harmonics $\widetilde{\psi}(k_\perp,\phi_k) = \sum_l \widetilde{\psi}_l(k_\perp) e^{il\phi_k}$. By Plancherel's identity,
\begin{equation}
  \braket{\Phi_{l,R}}{\psi} = \frac{i^{\abs{l}}}{\sqrt{2\pi}}\int_0^\infty k_\perp\,\dd k_\perp\,e^{-w^2 k_\perp^2/4}\widetilde{\psi}_l(k_\perp) J_{\abs{l}}(Rk_\perp).
\end{equation}
This is literally a Hankel transform! Because the Hankel transform preserves $L^2$ radial energy ($\int_0^\infty R\,\dd R\,\abs{\mathrm{H}_n[f](R)}^2 = \int_0^\infty k_\perp\,\dd k_\perp\,\abs{f(k_\perp)}^2$), integrating $\abs{\braket{\Phi_{l,R}}{\psi}}^2$ over $R\,\dd R$ and summing over $l$ yields:
\begin{equation}
  \norm{T\psi}^2 = \frac{1}{2\pi}\sum_{l\in\mathbb{Z}}\int_0^\infty k_\perp\,\dd k_\perp\,e^{-w^2 k_\perp^2/2}\abs{\widetilde{\psi}_l(k_\perp)}^2 = \int_{\mathbb{R}^2} e^{-w^2 k_\perp^2/2}\abs{\widetilde{\psi}(\mathbf{k})}^2\frac{\dd^2\mathbf{k}}{(2\pi)^2} \le \norm{\psi}_{L^2}^2.
\end{equation}
Reading off the operator $S_w$ in Fourier space: $\widetilde{(S_w\psi)}(\mathbf{k}) = e^{-w^2 k_\perp^2/2}\widetilde{\psi}(\mathbf{k})$! By the convolution theorem and Eq.~\eqref{eq:c2_gauss_shift}, multiplying by $e^{-w^2 k_\perp^2/2}$ in Fourier space is equivalent to convolving with a real-space Gaussian blur kernel $G_w(\mathbf{r}) = \frac{1}{2\pi w^2}e^{-r^2/(2w^2)}$!

\begin{paperhit}{JOSA A Eqs.~(7), (8), and (9) (p.~4)}
The continuous PVB frame operator is \emph{identically} a Gaussian low-pass blur of width $w$:
\[
  \widetilde{\Phi}_{l,R}(\mathbf{k}) = \sqrt{2\pi}\,(-i)^{\abs{l}} e^{il\phi_k} e^{-w^2 k_\perp^2/4} J_{\abs{l}}(Rk_\perp), \qquad
  (S_w\psi)(\mathbf{r}) = (G_w * \psi)(\mathbf{r}), \qquad \norm{T\psi}^2 \le 1\cdot\norm{\psi}_{L^2}^2.
\]
\end{paperhit}

\subsection{Step 4.3: Why Finite Aperture ($\mathrm{NA}$) Turns a Bessel Family Into a Sharp Frame}

Check the multiplier $e^{-w^2 k_\perp^2/2}$: as $k_\perp \to \infty$ (arbitrarily fine spatial ripples), $e^{-w^2 k_\perp^2/2} \to 0$! So on unconstrained $L^2(\mathbb{R}^2)$, there is NO positive lower bound $A > 0$. In any real optical instrument, however, a finite numerical aperture ($\mathrm{NA}$) cuts off spatial frequencies at $k_\perp \le k_{\max} = k_0\mathrm{NA}$ (the Paley--Wiener space $\mathcal{H}_{k_{\max}}$). On $[0, k_{\max}]$, the minimum of $e^{-w^2 k_\perp^2/2}$ is achieved at the edge $k_\perp = k_{\max}$ and its maximum at $k_\perp = 0$:

\begin{paperhit}{JOSA A Eq.~(10) and Table~1 (pp.~4--5)}
On the bandlimited optical space $\mathcal{H}_{k_{\max}}$, the PVB family is a genuine continuous frame with sharp bounds:
\[
  A(w,k_{\max}) = \exp\!\left(-\frac{w^2 k_{\max}^2}{2}\right), \qquad B = 1, \qquad \kappa(S_{w,k_{\max}}) = \frac{B}{A} = \exp\!\left(\frac{w^2 k_{\max}^2}{2}\right).
\]
Sharpness is verified by concentrating a test wavefield inside a thin Fourier annulus $k_{\max} - \epsilon \le k_\perp \le k_{\max}$ and letting $\epsilon \to 0^+$.
\end{paperhit}

% ==============================================================================
\section{Chapter 5: Propagating the PVB Frame Through Free Space}
\label{sec:c5}

Now combine Chapter~1 (the Fresnel propagator $\mathcal{U}_z$: multiply by $e^{-iz k_\perp^2/(2k_0)}$ in Fourier space) with Chapter~4 ($\widetilde{\Phi}_{l,R}$)!
Multiply $\widetilde{\Phi}_{l,R}(\mathbf{k})$ by $e^{-iz k_\perp^2/(2k_0)}$: the two $k_\perp^2$ exponents merge cleanly into $-\frac{1}{4}w_c^2(z)k_\perp^2$, where $w_c^2(z) \equiv w^2(1 + iz/z_R^{\mathrm{ring}})$ and $z_R^{\mathrm{ring}} = k_0 w^2/2$. Invert the 2D Fourier transform using Jacobi--Anger and Weber's integral~\eqref{eq:c2_weberJ} with complex $p = w_c^2(z)/4$ (valid for ALL $z \in \mathbb{R}$ because $\operatorname{Re} p = w^2/4 > 0$ strictly!):

\begin{paperhit}{JOSA A Eqs.~(11) and (12) (p.~5)}
Every PVB preserves its exact closed-form functional structure under free-space Fresnel propagation:
\[
  w_c^2(z) = w^2\!\left(1 + i\frac{z}{z_R^{\mathrm{ring}}}\right), \qquad
  \Phi_{l,R}^{(z)}(r,\phi) = \frac{1}{w_c^2(z)}\sqrt{\frac{2}{\pi}}\,e^{il\phi}\exp\!\left(-\frac{r^2+R^2}{w_c^2(z)}\right) I_{\abs{l}}\!\left(\frac{2Rr}{w_c^2(z)}\right).
\]
\end{paperhit}

Next, because $S_w$ (multiplier $e^{-w^2 k_\perp^2/2}$), $P_{k_{\max}}$ (indicator $\mathbf{1}_{k_\perp \le k_{\max}}$), and $\mathcal{U}_z$ (unit-modulus phase $e^{-iz k_\perp^2/(2k_0)}$) are ALL simple multiplications in Fourier space, they commute identically! Hence $S_w^{(z)} = \mathcal{U}_z S_w \mathcal{U}_z^\dagger = S_w$: the frame operator and its condition number are \textbf{completely invariant under propagation} \inpaper{Eq.~(13), p.~5}!

Finally, suppose an $\mathrm{LG}_{p,l}$ beam propagates to plane $z_m$, expanding by $\sigma_m = \sqrt{1+(z_m/z_R)^2}$ and acquiring wavefront curvature $e^{+ik_0 r^2/(2R_c(z_m))}$ plus Gouy phase $e^{-i(2p+\abs{l}+1)\vartheta_m}$. If we scale our detector's PVB template to $(R,w) \mapsto (\sigma_m R, \sigma_m w)$, notice that $\Phi_{l,\sigma_m R}(r,\phi;\,\sigma_m w) = \sigma_m^{-2}\Phi_{l,R}(r/\sigma_m,\phi;\,w)$ (because of the $1/w^2$ factor!). Multiplying by $\sigma_m^{+1} e^{+ik_0 r^2/[2R_c(z_m)]}$ cancels the curvature phase and makes the Jacobian $\mathbf{r} = \sigma_m\mathbf{r}'$ ($\dd^2\mathbf{r} = \sigma_m^2\dd^2\mathbf{r}'$) cancel $\sigma_m^{-1}\cdot\sigma_m^{-1}\cdot\sigma_m^2 = 1$ identically:

\begin{paperhit}{JOSA A Eqs.~(14) and (15) (p.~6)}
\[
  \phi_{l,R}^{(m)}(r,\phi) = \sigma_m e^{+ik_0 r^2/[2R_c(z_m)]}\Phi_{l,\sigma_m R}(r,\phi;\,\sigma_m w)
  \implies \braket{\phi_{l,R}^{(m)}}{\psi_{p,l}(z_m)} = e^{-i(2p+\abs{l}+1)\vartheta_m}\,c_{l,p}(R).
\]
\end{paperhit}

% ==============================================================================
\section{Chapter 6: Inverting the Gaussian Blur in the Presence of Noise}
\label{sec:c6}

\begin{whystep}[Why Naive Inversion Blows Up]
Suppose you measure the blurred frame output $\psi_{\mathrm{meas}} = S_w\psi + \varepsilon$, where $\varepsilon$ is tiny detector noise. If you naively divide by $g(k_\perp) = e^{-w^2 k_\perp^2/2}$ in Fourier space, noise at $wk_\perp = 3$ is multiplied by $e^{9/2} \approx 90$ in amplitude ($8100\times$ in variance!). How do we trade a tiny controlled bias for huge noise suppression?
\end{whystep}

Minimize the penalized objective $J[\psi] = \norm{S_w\psi - \psi_{\mathrm{meas}}}_{L^2}^2 + \alpha_{\mathrm{reg}}\norm{\psi}_{L^2}^2$ ($\alpha_{\mathrm{reg}} > 0$). Replace $\psi \to \psi + \epsilon\eta$ and set $\left.\frac{\dd}{\dd\epsilon}J[\psi+\epsilon\eta]\right|_{\epsilon=0} = 0$:
\begin{equation}
  2\operatorname{Re}\braket{\eta}{(S_w^2 + \alpha_{\mathrm{reg}}\mathbb{I})\psi - S_w\psi_{\mathrm{meas}}} = 0\quad\forall\eta
  \implies \hat{\psi}_{\alpha_{\mathrm{reg}}} = (S_w^2 + \alpha_{\mathrm{reg}}\mathbb{I})^{-1}S_w\psi_{\mathrm{meas}}.
\end{equation}
In Fourier space, the filter multiplier is $\widetilde{H}(g) = \frac{g}{g^2+\alpha_{\mathrm{reg}}} = (g + \alpha_{\mathrm{reg}}/g)^{-1}$. By the elementary AM--GM inequality ($g + \alpha_{\mathrm{reg}}/g \ge 2\sqrt{\alpha_{\mathrm{reg}}}$, with equality at $g = \sqrt{\alpha_{\mathrm{reg}}}$), the noise gain can NEVER exceed $(2\sqrt{\alpha_{\mathrm{reg}}})^{-1}$!
Writing $\hat{\psi}_{\alpha_{\mathrm{reg}}} - \psi = (S_{w,\alpha_{\mathrm{reg}}}^{-1}S_w - \mathbb{I})\psi + S_{w,\alpha_{\mathrm{reg}}}^{-1}\varepsilon$ splits the mean-squared error into $\mathrm{MSE} = \mathrm{Bias}^2 + \mathrm{Var}$. For white noise of power $P_\varepsilon$ on $k_\perp \le k_{\max}$ (density $\sigma_\varepsilon^2 = 4\pi P_\varepsilon/k_{\max}^2$), substituting $y = e^{-w^2 k_\perp^2}$ ($\dd y = -2w^2 k_\perp e^{-w^2 k_\perp^2}\dd k_\perp$) evaluates the variance integral in closed form:

\begin{paperhit}{JOSA A Eqs.~(16), (17), and (18) (pp.~6--7)}
\begin{align*}
  S_{w,\alpha_{\mathrm{reg}}}^{-1} &= (S_w^2 + \alpha_{\mathrm{reg}}\mathbb{I})^{-1}S_w, \qquad
  \widetilde{H}_{\max} = \frac{1}{2\sqrt{\alpha_{\mathrm{reg}}}}\text{ at } k_{\mathrm{peak}} = \frac{1}{w}\sqrt{\ln\!\left(\frac{1}{\alpha_{\mathrm{reg}}}\right)}, \\
  \mathrm{Bias}^2 &= \int_0^\infty\left[\frac{\alpha_{\mathrm{reg}}}{e^{-2au}+\alpha_{\mathrm{reg}}}\right]^2\rho_{2,1}(u)\,\dd u, \qquad
  \mathrm{Var} = \frac{P_\varepsilon}{(wk_{\max})^2}\left[\frac{1}{e^{-w^2 k_{\max}^2}+\alpha_{\mathrm{reg}}} - \frac{1}{1+\alpha_{\mathrm{reg}}}\right].
\end{align*}
\end{paperhit}

\begin{pencilbox}[Check the $258.65\times$ Deconvolution Gain by Hand ($wk_{\max} = 3.0$, $P_\varepsilon = 0.01$)]
At $\alpha_{\mathrm{reg}} = 0$: $\mathrm{MSE}(0) = \frac{0.01}{9}(e^9 - 1) = \mathbf{9.0023}$ (total disaster!). At $\alpha_{\mathrm{reg}}^{\mathrm{opt}} = 0.049025$: $\mathrm{Var} = \frac{0.01}{9}(\frac{1}{0.049148} - \frac{1}{1.049025}) = 0.02155$, $\mathrm{Bias}^2 = 0.01326$, $\mathrm{MSE}_{\min} = \mathbf{0.03481}$---a $\mathbf{258.65\times}$ ($\mathbf{24.13\,\mathrm{dB}}$) improvement!
\end{pencilbox}

% ==============================================================================
\section{Chapter 7: Connecting PVBs to Standard $\mathrm{LG}_{p,l}$ Modes and Finite Rings}
\label{sec:c7}

\subsection{Step 7.1: Exact $(p+1)$-Term Expansion $c_{l,p}(R) = \braket{\Phi_{l,R}}{\psi_{p,l}}$}

Every circular laser beam can be expanded in orthonormal Laguerre--Gaussian modes $\psi_{p,l}(r,\phi) = N_{p,l}(\frac{\sqrt{2}r}{w_0})^{\abs{l}}\allowbreak L_p^{\abs{l}}(\frac{2r^2}{w_0^2})\allowbreak e^{-r^2/w_0^2 + il\phi}$, where $L_p^{\abs{l}}(x) = \sum_{m=0}^p \frac{(-1)^m}{m!}\binom{p+\abs{l}}{p-m} x^m$. Setting $a_0 = \frac{1}{w^2} + \frac{1}{w_0^2}$ and $b_R = \frac{2R}{w^2}$ writes the overlap integral as \inpaper{Eq.~(25), p.~10}:
\begin{equation}
  c_{l,p}(R) = 2\pi C_0 N_{p,l} e^{-R^2/w^2}\left(\frac{\sqrt{2}}{w_0}\right)^{\abs{l}}\sum_{m=0}^p \frac{(-1)^m}{m!}\binom{p+\abs{l}}{p-m}\left(\frac{2}{w_0^2}\right)^m \int_0^\infty r^{\abs{l}+2m+1} e^{-a_0 r^2} I_{\abs{l}}(b_R r)\,\dd r.
\end{equation}
Insert the power series $I_{\abs{l}}(b_R r) = \sum_{k=0}^\infty \frac{(b_R r/2)^{\abs{l}+2k}}{k!\,(\abs{l}+k)!}$ and integrate term by term via $\int_0^\infty r^{2s+1}e^{-a_0 r^2}\dd r = \frac{s!}{2a_0^{s+1}}$ to obtain the confluent hypergeometric sum ${}_1F_1(\abs{l}+m+1;\,\abs{l}+1;\,\zeta_R)$ with $\zeta_R = \frac{b_R^2}{4a_0} = \frac{R^2 w_0^2}{w^2(w_0^2+w^2)}$ \inpaper{Eq.~(26), p.~10}! Then apply Kummer's identity ${}_1F_1(\abs{l}+m+1;\,\abs{l}+1;\,\zeta_R) = e^{\zeta_R}\,{}_1F_1(-m;\,\abs{l}+1;\,-\zeta_R)$ (proved from Euler's integral via $t = 1 - s$), which turns the infinite series into the finite polynomial $m!\,L_m^{\abs{l}}(-\zeta_R)$, cancelling $1/m!$ cleanly:

\begin{paperhit}{JOSA A Eqs.~(19) and (20) (p.~8)}
\[
  c_{l,p}(R) = \mathcal{A}_{p,l}(R)\sum_{m=0}^p (-1)^m\binom{p+\abs{l}}{p-m}\xi^m L_m^{\abs{l}}(-\zeta_R), \qquad
  \mathcal{A}_{p,l}(R) = 2^{\abs{l}/2+1}\sqrt{\frac{p!}{(p+\abs{l})!}}\,\frac{R^{\abs{l}} w_0^{\abs{l}+1}}{(w_0^2+w^2)^{\abs{l}+1}}e^{-R^2/(w_0^2+w^2)},
\]
with $\xi = \frac{2w^2}{w_0^2+w^2}$ and $\zeta_R = \frac{R^2 w_0^2}{w^2(w_0^2+w^2)}$.
\end{paperhit}

\subsection{Step 7.2: Radial Nodal Washout and the $0.573$ WKB Law}

An $\mathrm{LG}_{p,l}$ beam has $p$ bright concentric rings separated by dark nodal circles. When you pass it through the PVB frame operator $S_w$ (Gaussian blur of width $w$), how big can $w$ be before the blur washes out the rings? Measure fidelity $\mathcal{F} = \frac{\abs{\braket{\psi_{p,l}}{S_w\psi_{p,l}}}^2}{\norm{S_w\psi_{p,l}}^2} = \frac{(\int_0^\infty e^{-au}\rho_{p,l}(u)\dd u)^2}{\int_0^\infty e^{-2au}\rho_{p,l}(u)\dd u}$ ($a = w^2/w_0^2$, $u = w_0^2 k_\perp^2/2$).
\begin{itemize}[leftmargin=1.4em,itemsep=1pt,topsep=1pt]
  \item \textbf{Simplest case ($p = 0$):} $\rho_{0,l}(u) = \frac{u^{\abs{l}}e^{-u}}{\abs{l}!}$ gives $\mathcal{F}(0,l) = [1 - a^2/(1+a)^2]^{\abs{l}+1}$. Solving $\mathcal{F} = 0.95$ algebraically yields $w_{\mathrm{crit}}(0,l)$!
  \item \textbf{Many radial rings ($p \gg 1$):} By analogy with a 2D quantum harmonic oscillator of energy $E = 2N_p = 2p + \abs{l} + 1$, a classical particle spends time $\dd t \propto \dd k/v_k \propto \dd u/\sqrt{u(4N_p-u)}$ inside $u \in (0,4N_p)$. Substituting $u = 2N_p(1-\cos\theta)$ gives $\int_0^{4N_p} e^{-au}\rho_{\mathrm{WKB}}(u)\dd u = e^{-2N_p a} I_0(2N_p a)$, so $\mathcal{F}_{\mathrm{WKB}} = I_0(c)^2/I_0(2c) = 0.95 \implies c_* = 0.3277 \implies \sqrt{c_*} = \mathbf{0.5724}$!
\end{itemize}

\begin{paperhit}{JOSA A Eqs.~(21), (22), (23), and Table~2(a) (pp.~8--9)}
\[
  \mathcal{F}(p,l,w/w_0) = \frac{\bigl(\int_0^\infty e^{-au}\rho_{p,l}(u)\dd u\bigr)^2}{\int_0^\infty e^{-2au}\rho_{p,l}(u)\dd u}, \quad
  w_{\mathrm{crit}}(0,l) = w_0\left[(1 - \sqrt{1-q_l})^{-1}-1\right]^{1/2}, \quad
  w_{\mathrm{crit}}(p,l) \approx \frac{0.573\,w_0}{\sqrt{2p+\abs{l}+1}}.
\]
Hand check at $w = w_0, l = 1$: $\mathcal{F}(0,1,1) = \mathbf{9/16 = 0.56250}$, $\mathcal{F}(1,1,1) = \mathbf{81/256 = 0.31641}$.
\end{paperhit}

\subsection{Step 7.3: Discretizing to $N$ Concentric Rings (Szeg\H{o} Toeplitz Bound)}

In an experiment, we choose $N$ evenly spaced rings $R_n = R_{\min} + (n-1)\Delta R$. In the annular regime ($R_{\min} \gg w$), Step~4.1 shows the Gram matrix is $T_{nm} = q^{(n-m)^2}$ with $q = e^{-\Delta R^2/(2w^2)}$. Writing $q^{k^2} = \frac{1}{2\pi}\int_{-\pi}^\pi f(\theta) e^{ik\theta}\dd\theta$ with $f(\theta) = \sum_{k=-\infty}^\infty q^{k^2} e^{ik\theta} = \vartheta_3(\theta/2,q)$, any vector $\mathbf{v} \in \mathbb{C}^N$ has Rayleigh quotient $\frac{\mathbf{v}^\dagger T_N\mathbf{v}}{\norm{\mathbf{v}}^2} = \frac{\int_{-\pi}^\pi f(\theta)\abs{V(\theta)}^2\dd\theta}{\int_{-\pi}^\pi \abs{V(\theta)}^2\dd\theta}$, which is strictly trapped between $\min_\theta f(\theta) = f(\pi) = \vartheta_4(0,q)$ and $\max_\theta f(\theta) = f(0) = \vartheta_3(0,q)$!

\begin{paperhit}{JOSA A Eq.~(24) and Table~2(b) (p.~9)}
\[
  \kappa(T_N) < \kappa(T_\infty) = \frac{\vartheta_3(0,q)}{\vartheta_4(0,q)}, \qquad q = e^{-\Delta R^2/(2w^2)}, \qquad N_{\max} \approx \left\lfloor \frac{R_{\max}-R_{\min}}{2w}\right\rfloor + 1.
\]
At $\Delta R = 2w$ ($q = e^{-2} = 0.1353$): $\vartheta_3/\vartheta_4 = 1.2713/0.7300 = \mathbf{1.7416}$ (bounding every $N \ge 1$!).
\end{paperhit}

% ==============================================================================
\section*{Summary Map: Where Every JOSA A Equation Was Built From Scratch}
\addcontentsline{toc}{section}{Summary Map: Where Every JOSA A Equation Was Built From Scratch}

\begin{table}[ht!]
\centering
\scriptsize
\renewcommand{\arraystretch}{1.12}
\setlength{\tabcolsep}{5pt}
\begin{tabular}{@{}llll@{}}
\toprule
\textbf{Version 3 Chapter} & \textbf{Background Foundations Built} & \textbf{Paper Milestones Reached} & \textbf{JOSA A Pages} \\
\midrule
Ch.~1--2 (Warm-Up) & Maxwell $\to$ Fresnel, Gaussian beam, Jacobi--Anger, Weber, Hankel & \emph{Pure prerequisites (none assumed)} & --- \\
Ch.~3 (PVB Ring) & Ideal $\delta$-ring $\to$ $2f$ lens apodized Bessel--Gauss & \textbf{Eqs.~(1), (2), (3), (4)} & pp.~2--3 \\
Ch.~4 (Frames) & 3-vector Mercedes-Benz toy frame $\to$ continuous $S_w$ & \textbf{Eqs.~(5), (6), (7), (8), (9), (10)}, Table~1 & pp.~3--5 \\
Ch.~5 (Propagation) & Fresnel multiplier commutation & \textbf{Eqs.~(11), (12), (13), (14), (15)} & pp.~5--6 \\
Ch.~6 (Noise) & Variational calculus + AM--GM + Parseval & \textbf{Eqs.~(16), (17), (18)} & pp.~6--7 \\
Ch.~7 (LG \& Sampling) & Kummer Euler proof, 2D oscillator WKB, Szeg\H{o} Rayleigh quotient & \textbf{Eqs.~(19)--(26)}, Table~2 & pp.~8--10 \\
\bottomrule
\end{tabular}
\end{table}

\end{document}
